\documentclass[11pt]{article}
\usepackage[margin=1in]{geometry}
\usepackage{amsmath,amssymb,amsthm,booktabs,microtype}
\usepackage[colorlinks=true]{hyperref}
\newtheorem{theorem}{Theorem}[section]
\newtheorem{lemma}[theorem]{Lemma}
\newtheorem{corollary}[theorem]{Corollary}
\newcommand{\mathbbR}{\mathbb R}
\newcommand{\R}{\mathbb R}
\newcommand{\E}{\mathbb E}
\title{Learning narrow-network approximations\\with Gaussian Bayesian ReLU networks}
\author{ChatGPT-6 Astra \and Claude Opus 5.5 \and Kaarel H\"anni \and Dmitry Vaintrob\thanks{Authors are listed in alphabetical order.}}
\date{18 November 2026}
\hypersetup{pdftitle={Learning narrow-network approximations with Gaussian Bayesian ReLU networks},pdfauthor={ChatGPT-6 Astra; Claude Opus 5.5; Kaarel H\"anni; Dmitry Vaintrob}}
\begin{document}
\maketitle
\begin{abstract}
We give a Gaussian Bayesian neural-network learner an approximation guarantee measured only by the width and depth of a comparison network, without bounding its coefficients. The learner has one fixed dense architecture and an unrestricted independent Gaussian weight prior. A fixed final saturation is used in both the Gaussian likelihood and posterior prediction; a positive observation variance is selected from two empirical certificates. For bounded labels and an input law with density at most $\kappa$ relative to the uniform cube, the posterior mean competes with every width-$C$, hidden-depth-$h$ ReLU network, at every learner width $N\ge C+1$ and every prior exponent $\gamma\ge1$. At fixed input dimension $d$ and learner depth $H$, a sufficient sample count for additive error $\varepsilon$ is $\widetilde O_{d,H}(\kappa^{dp}(C+1)^{h(d+1)p}\varepsilon^{-(d+1)p})$, where $p=\max\{d+2,d(H+1)/2\}$. The learner depth must accommodate both the teacher and an $O(\log d)$ approximation construction. The proof combines a shifted multiscale interpolation of arbitrary-weight teachers with a Gaussian free-energy certificate and finite-parameter generalization. There is no prior over network sizes. The theorem places no restriction on Gaussian prior support.
\end{abstract}

\section*{The guarantee and its scope}
The best narrow network may have positive test loss, even when it captures useful structure in the data. An approximation oracle asks a Bayesian learner to approach that loss after enough observations, including when its own architecture is much wider. A width bound by itself allows arbitrarily large comparison coefficients; a useful statement must therefore account for sharp features without hiding a coefficient cap.

Theorem~\ref{thm:all-width-regular-Gaussian} supplies such an oracle under a regular input law. Its sample count is polynomial in comparison width at fixed dimension and depth, independent of learner width, and valid for noisy bounded responses. It applies to a Gaussian prior with independent coordinates throughout $\gamma\ge1$. The observation variance is calibrated from the training sample, and a fixed final saturation bounds each network output. These conventions are part of the learner.

The proof uses two different reasons that a posterior can generalize. At large widths, a network built by multiscale interpolation has a controlled Gaussian comparison cost. At smaller widths, ordinary parameter counting controls all bounded network outputs. The learner compares observable certificates for two positive temperatures and returns one posterior. The proof does not select the architecture or average over network sizes.

The dimensional dependence is substantial, and the result assumes a bounded-density input law. These hypotheses allow bounded-complexity approximations of sharp comparison functions. The theorem does not assert a coefficient-free oracle for arbitrary atomic input laws, nor does it concern an unbounded linear readout used directly in the likelihood.

\paragraph{Prior coefficient-reduction results.}
Coefficient reduction is not new to this analysis. Maass~\cite{maass} gives polynomial-bit, same-size representations for fan-out-one networks with Boolean outputs on finite rational input sets. Vardi and Shamir~\cite{vardi-shamir} approximate arbitrary-weight constant-depth ReLU networks by polynomial-weight, polynomial-width networks with a constant-factor depth increase under discrete and regular continuous input assumptions. The purpose here is to track the resources needed by the specified Bayesian prior and turn such approximations into an agnostic learning guarantee at every eligible learner width.

\section{Every learner width: a Bayesian model with a bounded output}
\label{sec:all-width-saturation}
Write $Q=[-1,1]^d$, let $U_Q$ be its uniform probability measure, and assume $|Y|\le1$ and $dP_X/dU_Q\le\kappa$. Squared population risk means $R(f)=\mathbb E(f(X)-Y)^2$. The class $\mathcal G_{C,h}$ consists of width-at-most-$C$, hidden-depth-$h$ ReLU networks with arbitrary finite real coefficients and biases and a linear readout. Put $H_d=2\lceil\log_2 d\rceil+2$.

The approximation theorem in Section~\ref{sec:width-only-regular} uses a wider network than the original comparison network. A second statistical certificate removes this width gap. It uses ordinary finite-parameter generalization at smaller widths, and the implicit-prior certificate at larger widths. Both certificates evaluate posteriors on the same fixed architecture; the learner selects one positive likelihood variance.

For this section, apply the fixed scalar activation
\[
 s(z)=\max\{-1,\min\{1,z\}\}
\]
after the linear readout. The Gaussian likelihood and the posterior prediction both use the bounded network output $s(z_w(x))$. Thus the prediction is the ordinary posterior mean of this network's output. This is a different observation model from a Gaussian likelihood on an unbounded readout followed only by clipping at prediction time.

The learner has $H$ hidden ReLU layers, each of width $N$, a scalar linear readout, and the final saturation. It has a constant input feature and no internal biases. Let $\pi_N$ be the fixed-architecture Gaussian prior below, with inverse scale $\omega_N$. Define
\[
 \widehat R_S(w)=\frac1n\sum_{i=1}^n(s(z_w(X_i))-Y_i)^2,
 \qquad
 \Pi_{S,\lambda}(dw)=
 \frac{e^{-\omega_N\widehat R_S(w)/(2\lambda)}}{Z_S(\lambda)}\pi_N(dw),
\]
so the physical Gaussian observation variance is $n\lambda/\omega_N$. The free energy is
\[
 F_S(\lambda)=-\frac{2\lambda}{\omega_N}\log Z_S(\lambda)
 =\inf_{\rho}\left\{\mathbb E_\rho\widehat R_S+
                   \frac{2\lambda}{\omega_N}\operatorname{KL}(\rho\Vert\pi_N)\right\}.
\]
The likelihood normalizer, which is independent of the network weights, is omitted in this definition. No prior over $N$ is used.

\subsection{A finite-parameter certificate at a positive temperature}
Let $\mathcal F_{N,H}$ be the class of saturated ReLU networks of hidden width $N$, hidden depth $H$, and the prescribed input features. Put
\[
 \widehat r_N(S)=\inf_{f\in\mathcal F_{N,H}}\frac1n\sum_i(f(X_i)-Y_i)^2.
\]
The parameter prior has every function in this class in their support on any finite input list. For the unrestricted Gaussian prior this follows from positive density at every finite parameter vector and continuity of the prediction map.

Consequently,
\[
 \lim_{\lambda\downarrow0}F_S(\lambda)=\widehat r_N(S).
\]
For example, the lower bound is immediate from the defining exponential integral. For the upper bound, any neighborhood with empirical loss at most $\widehat r_N+\eta$ has positive prior probability, so its contribution gives
$F_S(\lambda)\le\widehat r_N+\eta-2\lambda\log\pi_N(\text{neighborhood})/\omega_N$; first let $\lambda$ decrease and then let $\eta$ decrease.

Choose the first integer $j\ge0$ for which
\begin{equation}
 F_S(2^{-j})\le\widehat r_N(S)+1/n,
 \qquad \lambda_{\mathrm{par}}(S)=2^{-j}.
 \label{eq:parameter-temperature}
\end{equation}
This always exists and is strictly positive. Countable rational parameter approximations show that $\widehat r_N$ is measurable; continuity of the finite-data likelihood gives the same property for $F_S(2^{-j})$. The selected temperature is therefore measurable. This is an exact mathematical learning rule, with no computational-efficiency assertion.

A bounded real-valued class of pseudodimension $v$ obeys a relative uniform-convergence inequality for squared loss:
\begin{equation}
 R(f)\le\widehat R_S(f)+c\sqrt{\widehat R_S(f)a_N}+ca_N
 \quad\text{for every }f\in\mathcal F_{N,H},
 \qquad
 a_N=\frac{v\log(en)+\log(c(N+1)^2/\delta)}n,
 \label{eq:relative-parameter-bound}
\end{equation}
on one event of probability at least $1-\delta$ for all integer widths, for a numerical constant $c$. Allocate failure probability proportional to $\delta/(N+1)^2$ at width $N$. To obtain the real-loss form from the classical relative VC inequality~\cite{vapnik}, apply that inequality simultaneously to the superlevel sets of the bounded loss, then integrate their levels over $[0,4]$. The layer-cake identity gives the mean loss, and Cauchy--Schwarz bounds the integrated square-root term by twice the square root of empirical mean loss. The loss class has the same parameter-count bound, up to numerical or logarithmic factors, by adding a final subtraction and squaring gate.

For ReLU networks with saturation, if $D_N$ is their number of adjustable weights, one may take
\[
 v\le c D_N(H+2)\log(eD_N),\qquad
 D_N=N(d+2)+(H-1)N^2.
\]
Here the input has the original $d$ coordinates and a constant feature; changing its fixed normalization does not affect pseudodimension. These are the usual piecewise-linear neural-network parameter-count bounds~\cite{bartlettvc}. The final squaring in the loss changes only the fixed polynomial degree and a constant depth. In particular, at fixed $d,H$, $a_N=\widetilde O_{d,H}(N^2/n)$.

Integrate \eqref{eq:relative-parameter-bound} against any sample-dependent posterior and use concavity of the square root. Since its posterior empirical risk is at most its free energy, the observable quantity
\begin{equation}
 B_{\mathrm{par}}(S)=F_S(\lambda_{\mathrm{par}})
 +c\sqrt{F_S(\lambda_{\mathrm{par}})a_N}+ca_N
 \label{eq:parameter-certificate}
\end{equation}
is an upper bound on its posterior-average population risk. Jensen's inequality gives the same upper bound for its ordinary posterior mean.

\subsection{Selecting between two valid certificates}
Suppose the prior also has an implicit-complexity certificate: a measurable positive temperature $\lambda_{\mathrm{imp}}(S)$ and observable bound $B_{\mathrm{imp}}(S)$ on posterior-average risk, with a high-probability event on which this bound holds. The adaptive free-energy constructions supply exactly such a certificate, of the form $F_S(\lambda)+c_0\lambda$ plus a negligible tail term. For the saturated model they continue to hold: saturation is 1-Lipschitz, bounds outputs by one, does not increase comparison error to a bounded teacher, and only improves the moment and tail estimates.

Choose whichever of
\[
 (\lambda_{\mathrm{par}},B_{\mathrm{par}}),\qquad
 (\lambda_{\mathrm{imp}},B_{\mathrm{imp}})
\]
has smaller certificate, breaking a tie deterministically. The learner returns the ordinary mean of the single posterior at the selected temperature. On the intersection of the two certificate events its risk is at most
\begin{equation}
 \min\{B_{\mathrm{par}}(S),B_{\mathrm{imp}}(S)\}.
 \label{eq:minimum-certificates}
\end{equation}
This selection does not mix temperatures, functions, or architectures.

For a fixed comparison network $g$ representable in the width-$N$ class, Bernstein concentration of its bounded squared loss, \eqref{eq:parameter-temperature}, and \eqref{eq:parameter-certificate} give
\begin{equation}
 B_{\mathrm{par}}(S)
 \le R(g)+\widetilde O\left(\sqrt{R(g)D_N(H+2)/n}
                          +D_N(H+2)/n\right)
 \label{eq:parameter-comparison}
\end{equation}
on a further high-probability event. Here $g$ may itself be the saturated output of an arbitrary-weight ReLU network. Its original raw squared risk can replace $R(g)$ as an upper bound.

\subsection{The original independent Gaussian priors}
The two-certificate argument applies to the unrestricted Gaussian prior. For definiteness, write $D=d+1$, let the input feature be $\xi=(1,x)/\sqrt D$, and use independent centered coordinates with variances
\begin{equation}
 \operatorname{Var}(W^1_{ji})=N^{1-\gamma}/D,
 \qquad \operatorname{Var}(W^\ell_{ji})=N^{-\gamma}\ (2\le\ell\le H),
 \qquad \operatorname{Var}(v_j)=N^{-\gamma-1},
 \label{eq:width-only-Gaussian-prior}
\end{equation}
where $\gamma\ge1$ and $\omega_N=N^\gamma$. There is no restriction on the support of these weights. The weighted energy is
\[
 Q_N=\frac DN\|W^1\|_F^2+\sum_{\ell=2}^H\|W^\ell\|_F^2+N\|v\|_2^2.
\]

Here is the implementation bound needed to use the Gaussian comparison theorem. A width-$m$ ReLU network of path variation $P$ can be rescaled, without changing its function, so that every internal absolute spectral norm is at most one and the first and last squared Frobenius norms are both $P$. The path-flow proof is as follows. With forward masses $a^1_j=\|W^1_{j,:}\|_2$ and $a^\ell=|W^\ell|a^{\ell-1}$, and backward masses $b^H=|v|$, $b^{\ell-1}=|W^\ell|^{\mathsf T}b^\ell$, every layer has $\sum_j a_jb_j=P$. Remove neurons on no nonzero path and rescale neuron $j$ by $\sqrt{b_j/a_j}$. The transformed nonnegative internal matrices send $r^{\ell-1}$ to $r^\ell$ and their transposes send $r^\ell$ back, where $r_j=\sqrt{a_jb_j}$. Weighted Cauchy--Schwarz gives operator norm at most one. The endpoint squared norms are $\sum_jr_j^2=P$.

Zero-pad this implementation to width $N\ge m$. Multiply every hidden activation by $\sqrt N/D^{1/4}$, which rescales only the first matrix and readout. Its weighted energy is then bounded by
\begin{equation}
 E\le2\sqrt D\,P+(H-1)m.
 \label{eq:path-to-Gaussian-energy}
\end{equation}
The internal bound uses squared Frobenius norm at most rank times squared operator norm. These support properties are used to construct one comparison implementation; the Gaussian learner's prior remains unrestricted.

The selective Gaussian comparison in Section~\ref{sec:selective-standalone} shrinks only incoming rows of active teacher neurons, and the readout. Its normalized entropy cost is
\begin{equation}
 A=\widetilde O_{d,H}(1+P+m),
 \label{eq:Gaussian-path-resource}
\end{equation}
using \eqref{eq:path-to-Gaussian-energy}. The logarithmic factors include $n$, the depth, and this energy. This is the all-$\gamma\ge1$ selective comparison, so shrinking an $N$-dimensional readout or an active incoming row does not introduce an uncontrolled $N$ factor after division by $\omega_N=N^\gamma$.

The finite-dimensional saturated-output Gaussian certificate, proved in Section~\ref{sec:saturated-specialization}, has
\begin{equation}
 p=\max\{d+2,d(H+1)/2\},\qquad \rho=(d+2)/2,
 \label{eq:saturated-exponents}
\end{equation}
and, for a bounded teacher of risk $R_*$, gives
\begin{equation}
 B_{\mathrm{imp}}\le R_*+
 \widetilde O_{d,H}\left(
       A R_*^{1-\rho/p}n^{-1/p}+A^{p/\rho}n^{-1/\rho}\right).
 \label{eq:saturated-Gaussian-certificate}
\end{equation}
In the operator-norm cover proof, the cover logarithm has size $k^{d(H+1)/2}a^{-d}$, while saturation makes the loss envelope independent of $k$. The common free-energy root is
$\lambda^p=\theta_n(2F_S(\lambda)+\lambda b_n)^{p-\rho}$.
The same finite-cover and weighted-tail argument therefore gives \eqref{eq:saturated-Gaussian-certificate}; it is precisely the bounded-output specialization of the all-$\gamma$ Gaussian oracle. For bounded $R_*$, $n=\widetilde O_{d,H}((A/\varepsilon)^p)$ is sufficient for excess $\varepsilon$.

\begin{theorem}[Width-only oracle at every eligible learner width]
\label{thm:all-width-regular-Gaussian}
Assume $|Y|\le1$ and $dP_X/dU_Q\le\kappa$. Fix $C,h\ge1$ and a learner hidden depth $H\ge\max\{h,H_d\}$. Use the independent Gaussian prior \eqref{eq:width-only-Gaussian-prior}, final saturation in both likelihood and prediction, and the two-certificate calibration above. Then, on one event of probability at least $1-\delta$, simultaneously for every $\gamma\ge1$ and every $N\ge C+1$, the ordinary posterior mean satisfies
\[
 R(\widehat f_n)\le\inf_{g\in\mathcal G_{C,h}}R(g)+\varepsilon
\]
whenever
\begin{equation}
 n\ge\widetilde O_{d,H}\left(
   \kappa^{dp}(C+1)^{h(d+1)p}\varepsilon^{-(d+1)p}\right),
 \qquad p=\max\{d+2,d(H+1)/2\}.
 \label{eq:all-width-regular-sample}
\end{equation}
The comparison class has no coefficient bound. The extra neuron in $N\ge C+1$ carries the constant needed to compile arbitrary internal biases. If that constant wire is already included in the stated teacher width, $N\ge C$ suffices. The prior and calibration use none of $C,h,\kappa$, or the comparator's coefficients.
\end{theorem}
\begin{proof}
Choose a deterministic comparator within $\varepsilon/8$ of the risk infimum. Apply the hierarchical approximation theorem with a dyadic mesh $\Delta$ within a factor two below $\varepsilon/[64\kappa d(C+1)^h]$. Its interpolant $p_g$ has risk at most $R(g)+\varepsilon/8$, width
\[
 m=O_d(\Delta^{-d}),
\]
and path variation $P=O_d((C+1)^h\Delta^{-d})$. Pad by identity layers to depth $H$. Thus \eqref{eq:Gaussian-path-resource} is bounded by
\[
 A=\widetilde O_{d,H}\left(
      \kappa^d(C+1)^{h(d+1)}\varepsilon^{-d}\right).
\]
If $N\ge m$, the implicit certificate applied to this fixed interpolant is at most $R(g)+\varepsilon/2$ under \eqref{eq:all-width-regular-sample}.

If $C+1\le N<m$, compile the original comparator itself, including its biases, at width $N$ and depth $H$. Its saturated output belongs to $\mathcal F_{N,H}$ and has risk no larger than the original raw comparator. The parameter certificate gives excess at most $\varepsilon/2$ when
\[
 n\ge\widetilde O_{d,H}(m^2/\varepsilon^2)
 =\widetilde O_{d,H}(\kappa^{2d}(C+1)^{2dh}\varepsilon^{-(2d+2)}).
\]
This is implied by \eqref{eq:all-width-regular-sample}, because $p\ge d+2\ge2$, $\kappa\ge1$, $C+1\ge1$, and $\varepsilon\le1$. The selected posterior obeys the smaller certificate in either case. Allocate confidence among the uniform certificates and the two fixed-comparator concentration bounds; their union gives the claim. The parameter certificate uses a summable allocation over learner widths. Its additional logarithm of $N$ is bounded by a logarithm of $m$ in the only branch where that certificate is needed, and therefore introduces no learner-width dependence in the final sample bound.
\end{proof}

At fixed $d,H,\kappa$, the theorem competes with width $n^a$ for any $a<1/[h(d+1)p]$, with a vanishing additive error. The distribution may have arbitrary bounded label noise and the best comparison risk need not vanish. The restrictions doing work here are regularity of the input law and the fixed bounded output activation, not a coefficient bound or a prior on architecture sizes.


\section{Competing with arbitrary-weight networks under a regular input law}
\label{sec:width-only-regular}
The path-variation oracle also yields a width-only approximation guarantee under a restriction on the input distribution. The comparison network may have arbitrarily large coefficients. The learner still uses a single fixed architecture, but its width must be large enough for a bounded-variation approximation of the comparison function.

Let $Q=[-1,1]^d$, and let $U_Q$ denote its uniform probability measure. Assume $|Y|\le1$ and
\[
 \frac{dP_X}{dU_Q}\le\kappa,\qquad \kappa\ge1.
\]
The learner receives the augmented input $\xi(x)=(1,x)/\sqrt{d+1}$, which lies in the Euclidean unit ball. Its mixed path variation therefore uses first-row Euclidean norms in these augmented coordinates. Put
\[
 H_d=2\lceil\log_2 d\rceil+2,\qquad
 J_d=10d^2\sqrt{d+1}.
\]

\begin{theorem}[A bounded-path approximation with no coefficient bound]
\label{thm:width-only-approximation}
Let $g:\mathbb R^d\to\mathbb R$ be computed by a ReLU network with $h\ge1$ hidden layers, at most $C\ge1$ neurons in each hidden layer, arbitrary finite real weights and biases, and a linear readout. Define
\[
 R_0=(C+1)^h,\qquad f=\max\{-1,\min\{1,g\}\}.
\]
For every $0<\Delta\le1$, there is a ReLU network $p$ with no internal biases, augmented input $\xi$, hidden depth $H_d$, and
\begin{align}
 \operatorname{width}(p)&\le(4d+1)(5/\Delta)^d,\label{eq:width-only-width}\\
 P(p)&\le J_d5^d\Delta^{-(d+1)},\label{eq:width-only-path}\\
 |p(x)|&\le1\quad(x\in Q),\qquad
 \int_Q|p-f|\,dU_Q\le dR_0\Delta.\label{eq:width-only-error}
\end{align}
Consequently,
\begin{equation}
 R(p)\le R(g)+4\kappa dR_0\Delta.
 \label{eq:width-only-risk}
\end{equation}
Identity layers increase the hidden depth to any $H\ge H_d$ without changing the width or path variation bounds.
\end{theorem}

\subsection{Variation on coordinate lines}
Fix every input coordinate except one. The restricted ReLU network is a univariate continuous piecewise affine function. At the first hidden layer, the breakpoints of all $C$ neurons divide the line into at most $C+1$ intervals. On each interval on which all preceding features are affine, the next layer creates at most $C$ new breakpoints. Inductively the output has at most $(C+1)^h=R_0$ affine intervals. Biases and arbitrary coefficient magnitudes do not change this count.

On each affine interval, clipping the scalar output to $[-1,1]$ gives a monotone function with total variation at most two. Thus on every coordinate line,
\begin{equation}
 \operatorname{TV}(f\text{ on that line})\le2R_0.
 \label{eq:width-only-line-TV}
\end{equation}
This holds on the whole line, including lines outside $Q$.

The following translation estimate is useful. For any $z\in\mathbb R^d$,
\begin{equation}
 \int_Q|f(x+z)-f(x)|\,dU_Q(x)\le R_0\|z\|_1.
 \label{eq:width-only-translation}
\end{equation}
Indeed, telescope one coordinate at a time. For a one-dimensional function of total variation at most $2R_0$, its integral translation error on an interval is at most $2R_0|z_j|$. Integrating the other coordinates over intervals of length two and dividing by $2^d$ proves the asserted contribution $R_0|z_j|$. The variation estimate is uniform over those other coordinates and their intermediate translations.

\subsection{Kuhn hats and a randomly shifted grid}
Define
\[
 \Phi(z)=\operatorname{ReLU}\!\left(1-
       \max\{0,z_1,\ldots,z_d\}
       -\max\{0,-z_1,\ldots,-z_d\}\right).
\]
For a grid vertex $v\in u+\Delta\mathbb Z^d$, set
$\phi_v(x)=\Phi((x-v)/\Delta)$. These are the nodal hats of the Kuhn triangulation. They are nonnegative, vanish unless $\|x-v\|_\infty\le\Delta$, and satisfy
\begin{equation}
 \sum_{v\in u+\Delta\mathbb Z^d}\phi_v(x)=1.
 \label{eq:width-only-partition}
\end{equation}
Here is a direct verification. In a unit grid cell, order the fractional coordinates as $\theta_{\pi(1)}\ge\cdots\ge\theta_{\pi(d)}$. The only potentially nonzero hats correspond to the $d+1$ vertices obtained by adding the first $j$ ordered coordinate vectors to the lower corner. Their values are
\[
 1-\theta_{\pi(1)},\quad
 \theta_{\pi(1)}-\theta_{\pi(2)},\quad\ldots,\quad
 \theta_{\pi(d-1)}-\theta_{\pi(d)},\quad
 \theta_{\pi(d)}.
\]
They sum to one. A different vertex either lies outside the adjacent grid cube or has a selected coordinate with smaller fractional part than an unselected one; the displayed maximum-plus-maximum formula then gives zero. Ties and cell boundaries follow by continuity. Integrating \eqref{eq:width-only-partition} over a fundamental grid cell also gives $\int_{\mathbb R^d}\Phi=1$.

For a shift $u\in[0,\Delta)^d$, form the nodal interpolant
\[
 I_uf(x)=\sum_{v\in u+\Delta\mathbb Z^d}f(v)\phi_v(x).
\]
It stays in $[-1,1]$. Fixed grid vertices can meet narrow features of a function, so a deterministic unshifted grid is not sufficient for the next estimate. Averaging over a uniformly random shift removes that issue. By \eqref{eq:width-only-partition}, Tonelli's theorem, and the tiling of $\mathbb R^d$ by the translated fundamental cells,
\begin{align*}
 \mathbb E_u\int_Q|I_uf-f|\,dU_Q
 &\le\int_{\mathbb R^d}\Phi(-z)
       \int_Q|f(x+\Delta z)-f(x)|\,dU_Q(x)\,dz\\
 &\le R_0\Delta\int\Phi(-z)\|z\|_1\,dz
 \le dR_0\Delta.
\end{align*}
The last inequality uses $\operatorname{supp}\Phi\subset[-1,1]^d$ and $\int\Phi=1$. Some fixed shift therefore has the claimed error. Only vertices in $[-1-\Delta,1+\Delta]^d$ can contribute on $Q$, and there are at most
\[
 M\le(2/\Delta+3)^d\le(5/\Delta)^d
\]
 such vertices. Retain only these hats.

\subsection{Realizing the interpolant and bounding its path variation}
Each hat has a ReLU realization with $H_d$ hidden layers. In the first layer, compute
\[
 \operatorname{ReLU}((x_j-v_j)/\Delta),\quad
 \operatorname{ReLU}((v_j-x_j)/\Delta)\quad(1\le j\le d),
 \qquad \operatorname{ReLU}(1)=1.
\]
The last neuron carries a constant through subsequent layers. Compute the maxima of the positive and negative lists by balanced binary trees, padding with zeros to $D=2^{\lceil\log_2d\rceil}$ leaves. For nonnegative $a,b$,
\[
 \max\{a,b\}=\operatorname{ReLU}(\operatorname{ReLU}(a-b)+b).
\]
Thus each tree level takes two ReLU layers. The final layer computes $\operatorname{ReLU}(1-M_+-M_-)$. It uses the carried constant, so no internal bias is needed. Every layer has at most $4d+1$ neurons per hat.

For completeness, path budgets are propagated by absolute weights. At a first-layer neuron, its budget is its incoming Euclidean row norm. At a later neuron it is the sum of its predecessors' budgets times the corresponding absolute weights. The output budget is exactly mixed path variation. In augmented coordinates, each first-layer hinge has budget at most
\[
 \sqrt{d+1}\sqrt{1+v_j^2}/\Delta\le\sqrt{5(d+1)}/\Delta,
\]
while the constant has budget $\sqrt{d+1}$. A two-layer maximum block sends budgets $p_a,p_b$ to $p_a+2p_b\le2(p_a+p_b)$. Hence either maximum tree has budget at most
\[
 Dd\sqrt{5(d+1)}/\Delta\le2d^2\sqrt{5(d+1)}/\Delta.
\]
The final hat therefore has budget at most
\[
 \sqrt{d+1}+4d^2\sqrt{5(d+1)}/\Delta
 \le J_d/\Delta.
\]
Place the $M$ hat networks side by side and use readout coefficients $f(v)\in[-1,1]$. This proves \eqref{eq:width-only-width} and \eqref{eq:width-only-path}. Identity padding preserves these budgets because every hidden coordinate is nonnegative.

Finally clipping $g$ can only reduce squared loss when $|Y|\le1$. Both $p$ and $f$ are in $[-1,1]$, so
\[
 |(p-Y)^2-(f-Y)^2|\le4|p-f|.
\]
Density domination and the interpolation estimate prove \eqref{eq:width-only-risk}, finishing Theorem~\ref{thm:width-only-approximation}.

\subsection{A sharper path bound from hierarchical interpolation}
A multiscale representation reduces the path variation further. Define
\[
 K_d=J_d\bigl(5^d+2d4^{d-1}\bigr).
\]

\begin{theorem}[Hierarchical bounded-path approximation]
\label{thm:hierarchical-width-only}
Under the assumptions of Theorem~\ref{thm:width-only-approximation}, let $\Delta=2^{-J}$ for an integer $J\ge0$. There is a ReLU network $p$ of hidden depth $H_d$ such that
\begin{align}
 \operatorname{width}(p)&\le2(4d+1)5^d\Delta^{-d},\notag\\
 P(p)&\le2K_d R_0\Delta^{-d},\label{eq:hierarchical-path}\\
 |p(x)|&\le1\quad(x\in Q),\qquad
 \int_Q|p-f|\,dU_Q\le2dR_0\Delta.\notag
\end{align}
In particular $R(p)\le R(g)+8\kappa dR_0\Delta$.
\end{theorem}
\begin{proof}
Choose one shift $u$ uniformly in $[0,1)^d$ and use all nested grids $u+\Delta_j\mathbb Z^d$, where $\Delta_j=2^{-j}$ for $0\le j\le J$. Write $I_jf$ for the corresponding Kuhn interpolant.

The triangulations are nested. Indeed, a Kuhn triangulation of mesh $\Delta_j$ is cut out by the hyperplanes $x_i-u_i=k\Delta_j$ and $x_i-x_l-(u_i-u_l)=k\Delta_j$. Every coarse hyperplane is a fine hyperplane. Consequently every function affine on the coarse simplices is reproduced by fine interpolation. We therefore have the exact telescoping representation
\begin{equation}
 I_Jf=I_0f+\sum_{j=1}^J\sum_{v\in u+\Delta_j\mathbb Z^d}
       c_{j,v}\phi_{j,v},\qquad
 c_{j,v}=f(v)-I_{j-1}f(v).
 \label{eq:hierarchical-expansion}
\end{equation}
Only hats whose supports meet $Q$ need be retained.

There is a simple formula for the detail coefficients. Write a fine vertex uniquely as
$v=u+\Delta_jb+2\Delta_j k$, with $b\in\{0,1\}^d$ and $k\in\mathbb Z^d$. Its fractional coordinates in the containing coarse cell are $b/2$. The barycentric formula above then gives
\begin{equation}
 I_{j-1}f(v)=\tfrac12\bigl[f(v-\Delta_j b)+f(v+\Delta_j b)\bigr].
 \label{eq:hierarchical-midpoint}
\end{equation}
This also holds for $b=0$, when the detail is zero.

Put $Q_j=[-1-\Delta_j,1+\Delta_j]^d$. Averaging over $u$ is equivalent to averaging modulo the coarse mesh $2\Delta_j$. Tiling separately for each parity $b$ and using \eqref{eq:hierarchical-midpoint} gives
\begin{align*}
 \mathbb E_u\sum_{v\in (u+\Delta_j\mathbb Z^d)\cap Q_j}|c_{j,v}|
 &\le(2\Delta_j)^{-d}\sum_{b\in\{0,1\}^d}
 \frac12\int_{Q_j}\bigl(
 |f(v)-f(v-\Delta_j b)|+|f(v)-f(v+\Delta_j b)|\bigr)\,dv\\
 &\le dR_0(2+2\Delta_j)^{d-1}\Delta_j^{1-d}
 \le d4^{d-1}R_0\Delta_j^{1-d}.
\end{align*}
For the second inequality, the unnormalized coordinate translation bound on a box of side $2+2\Delta_j$ is
$2R_0(2+2\Delta_j)^{d-1}\|z\|_1$, and
$\sum_{b\in\{0,1\}^d}\|b\|_1=d2^{d-1}$.

Implement \eqref{eq:hierarchical-expansion} by placing all of its hats in parallel. A hat of mesh $\Delta_j$ has path budget at most $J_d/\Delta_j$, as already proved. The base interpolant costs at most $J_d5^d$. The expected sum of path budgets is consequently at most
\[
 J_d5^d+J_dd4^{d-1}R_0\sum_{j=1}^J\Delta_j^{-d}
 \le K_dR_0\Delta^{-d}.
\]
The total number of retained hats is at most
$5^d\sum_{j=0}^J\Delta_j^{-d}\le2\cdot5^d\Delta^{-d}$, which gives the stated width. Every hat has depth $H_d$.

The final interpolant $I_Jf$ still takes values in $[-1,1]$, despite the signed multiscale coefficients. The earlier random-shift proof gives
$\mathbb E_u\int_Q|I_Jf-f|\,dU_Q\le dR_0\Delta$.
The sum of this error divided by $dR_0\Delta$ and the constructed path budget divided by $K_dR_0\Delta^{-d}$ has expectation at most two. Some shift makes the sum at most two. For that shift, each nonnegative term is at most two, proving both claims simultaneously. Density domination and the bounded-output squared-loss inequality give the final risk estimate.
\end{proof}


\section{Hyperplane anti-concentration in place of a density bound}
\label{sec:hyperplane-regularity}
The interpolation argument uses regularity relative to volume. A different compiler uses only anti-concentration near affine hyperplanes. This includes some distributions supported on a lower-dimensional surface, and gives a different tradeoff between comparison width and approximation accuracy.

Assume $X\in Q=[-1,1]^d$ and, for some $\kappa\ge1$ and $0<\alpha\le1$,
\begin{equation}
 \Pr\{|u^{\mathsf T}X-t|\le s\}\le\kappa s^\alpha
 \quad\text{for every }\|u\|_2=1,\ t\in\mathbb R,\ s>0.
 \label{eq:hyperplane-anticoncentration}
\end{equation}
In particular every affine hyperplane has probability zero. A bounded density relative to the uniform cube implies this condition with $\alpha=1$ and slab constant at most $\sqrt d$ times its density bound: choose a coordinate of the unit normal of magnitude at least $1/\sqrt d$ and integrate that coordinate first. The uniform law on the unit sphere in $\mathbb R^d$, $d\ge3$, also satisfies \eqref{eq:hyperplane-anticoncentration} with $\alpha=1$ and a constant of order $\sqrt d$. For the circle in $\mathbb R^2$, exponent $\alpha=1/2$ suffices. These follow directly from the one-coordinate marginal densities of the sphere. No density in the ambient cube is required.

\begin{theorem}[A compiler under hyperplane anti-concentration]
\label{thm:hyperplane-compiler}
Let $g$ be a width-$C$, hidden-depth-$h$ ReLU network with arbitrary coefficients and biases. Put
\[
 M=(C+1)^{dh},\qquad T=M(Ch+1).
\]
Under \eqref{eq:hyperplane-anticoncentration}, for every $0<\varepsilon\le1$ there is a four-hidden-layer ReLU network $p$, with augmented input $\xi=(1,x)/\sqrt{d+1}$ and no internal biases, such that
\begin{align}
 \operatorname{width}(p)&\le3T,\notag\\
 |p(x)|&\le1\quad(x\in Q),\notag\\
 R(p)&\le R(g)+\varepsilon/4,\label{eq:hyperplane-compiler-risk}\\
 P(p)&\le80(d+1)64^{1/\alpha}
       \kappa^{1/\alpha}T^{1+1/\alpha}\varepsilon^{-1/\alpha}.
 \label{eq:hyperplane-compiler-path}
\end{align}
Identity padding gives the same statement at every hidden depth $H\ge4$.
\end{theorem}
\begin{proof}
A ReLU network is affine on each of its full-dimensional activation cells. At one layer, $C$ hyperplanes cut any preceding convex cell into at most $\sum_{j=0}^d\binom Cj\le(C+1)^d$ cells. Induction gives at most $M$ full-dimensional cells. Each is described by at most $Ch$ affine inequalities, one for each hidden activation along its activation pattern. Assign an identically zero preactivation one fixed activation convention, so that the cells form a disjoint polyhedral partition rather than duplicated sign patterns. Degenerate identically zero constraints can then be removed. The complement of the cell interiors is contained in a finite union of affine hyperplanes and has probability zero by \eqref{eq:hyperplane-anticoncentration}.

Write $R_Q=\sqrt d$ and $D=d+1$. Normalize the nonconstant inequalities defining a cell as
$u_j^{\mathsf T}x-t_j\ge0$ with $\|u_j\|_2=1$. If $t_j<-R_Q$, the inequality is strict everywhere on $Q$ and can be removed. If $t_j>R_Q$, the cell does not meet $Q$ and can be discarded. All retained offsets satisfy $|t_j|\le R_Q$. Let $r$ be the number retained in a particular cell, so $r\le Ch$.

For a scale $0<\delta\le1$, define
\[
 a_j(x)=\operatorname{ReLU}((u_j^{\mathsf T}x-t_j)/\delta)
       -\operatorname{ReLU}((u_j^{\mathsf T}x-t_j)/\delta-1),
 \qquad
 b(x)=\operatorname{ReLU}\left(\sum_{j=1}^r a_j(x)-(r-1)\right).
\]
Every $a_j$ lies in $[0,1]$, as does $b$. The gate $b$ is zero outside the cell interior, equals one when all its retained inequalities have slack at least $\delta$, and is one identically when $r=0$. Gates from different cell interiors cannot be positive at the same point. Values on a shared facet are zero, so this disjointness also ensures the boundedness assertion below on all of $Q$.

On the cell, the teacher output has the affine form $\ell(x)=a^{\mathsf T}x+b_0$. We next replace its saturation by a saturated affine function with a controlled slope. Fix $B\ge1$. If $\|a\|_2\le B$, retain $\ell$ unless it already has constant saturated sign on $Q$, in which case use that constant. Its retained bias can be bounded by $BR_Q+1$: a larger absolute bias would force a constant saturated sign. If $\|a\|_2>B$, write $\ell=\|a\|_2(u^{\mathsf T}x-t)$ and replace its slope by $B$. If $|t|>R_Q+1/B$, both saturated functions are the same constant on $Q$; otherwise the new bias is at most $BR_Q+1$ in absolute value. Denote the resulting saturated affine output by $\phi(x)\in[-1,1]$. Its coefficients have slope norm at most $B$ and bias magnitude at most $BR_Q+1$. It differs from $s(\ell(x))$ only in a slab of width $1/B$, and its absolute difference there is at most two.

Combine the region output with its gate by the exact ReLU mask
\[
 m(x)=\operatorname{ReLU}(\phi(x)+b(x)-1)
      -\operatorname{ReLU}(-\phi(x)+b(x)-1).
\]
It equals $\phi$ when $b=1$, is zero when $b=0$, and has absolute value at most $b$. Sum these masks over the retained cells to obtain $p$. At most one gate is positive on $Q$, so $|p|\le1$.

Except on the union of the facet slabs of width $\delta$ and the output slabs of width $1/B$, the resulting function equals $s(g)$. There are at most $MCh$ of the former and at most $M$ of the latter. Thus
\[
 \mathbb E|p(X)-s(g(X))|
 \le2\kappa\left[MCh\,\delta^\alpha+MB^{-\alpha}\right].
\]
Take
\[
 \delta=\left(\frac{\varepsilon}{64\kappa T}\right)^{1/\alpha},
 \qquad B=\left(\frac{64\kappa M}{\varepsilon}\right)^{1/\alpha}.
\]
The last display is at most $\varepsilon/16$. Squared loss between bounded predictions and labels changes by at most four times prediction difference, and clipping the original teacher cannot hurt its risk. This proves \eqref{eq:hyperplane-compiler-risk}.

Here is an explicit depth and path count. At the first hidden layer, compute the two hinges of each $a_j$, the two hinges representing
\[
 \phi=\operatorname{ReLU}(\ell'+1)-\operatorname{ReLU}(\ell'-1)-1,
\]
and one constant feature. At the second layer compute the $a_j$ and the positive and negative parts of $\phi$, carrying the constant. At the third layer compute $b$ and pass those parts and the constant. The fourth layer computes the two mask hinges. This uses at most $\max\{4,2r+3\}\le3(Ch+1)$ neurons per cell and no internal biases.

In the augmented input coordinates, the first row norm of a facet hinge is at most $3D/\delta$, and that of an output hinge is at most $4DB$. The constant row has norm $\sqrt D$. Absolute path budgets therefore give at most $6D/\delta$ for $a_j$, at most $9DB$ for either part of $\phi$, and at most $7D(Ch+1)/\delta$ for the gate. Each of the two final mask hinges has budget at most $19DB+7D(Ch+1)/\delta$. The cell's total contribution is at most
\[
 40D\left(B+(Ch+1)/\delta\right).
\]
Since $B\le1/\delta$, summing the $M$ cells gives at most
$80DT/\delta$, which is \eqref{eq:hyperplane-compiler-path}. This also proves the claimed width and depth.
\end{proof}

\begin{corollary}[An all-width Gaussian oracle under slab regularity]
\label{cor:hyperplane-oracle}
Assume $|Y|\le1$ and \eqref{eq:hyperplane-anticoncentration}. Fix a comparison width $C$, comparison hidden depth $h$, and learner hidden depth $H\ge\max\{h,4\}$. The same unrestricted Gaussian prior, final saturation, and two-certificate learner satisfy, simultaneously for all $\gamma\ge1$ and $N\ge C+1$,
\[
 R(\widehat f_n)\le\inf_{g\in\mathcal G_{C,h}}R(g)+\varepsilon
\]
with probability at least $1-\delta$ as soon as
\begin{equation}
 n\ge\widetilde O_{d,H,\alpha}\left(
 \kappa^{p/\alpha}
 \bigl[(Ch+1)(C+1)^{dh}\bigr]^{p(1+1/\alpha)}
 \varepsilon^{-p(1+1/\alpha)}\right),
 \qquad p=\max\{d+2,d(H+1)/2\}.
 \label{eq:hyperplane-sample}
\end{equation}
\end{corollary}
\begin{proof}
The Gaussian path-to-energy comparison applies to the compiler of Theorem~\ref{thm:hyperplane-compiler}, whose width is $m\le3T$ and whose resource is of order
$\kappa^{1/\alpha}T^{1+1/\alpha}\varepsilon^{-1/\alpha}$, up to logarithms. At $N\ge m$, the Gaussian implicit certificate therefore gives the claimed sample bound. At $C+1\le N<m$, the parameter certificate needs only $\widetilde O_{d,H}(T^2/\varepsilon^2)$ samples. This is dominated by \eqref{eq:hyperplane-sample}, since $p\ge2$, $0<\alpha\le1$, and $\kappa,T\ge1$. Select a fixed comparator arbitrarily close to the infimum and allocate confidence exactly as in Theorem~\ref{thm:all-width-regular-Gaussian}.
\end{proof}

For bounded-density laws, one may use either compiler. The multiscale compiler gives a better power of comparison width, while this polyhedral compiler gives a better power of approximation accuracy when $d>1$ and $\alpha=1$. The latter also needs only four hidden layers for its comparison implementation. The actual calibration rule does not need to know which argument is applicable.

\section{A smooth bounded output: final tanh}
\label{sec:smooth-output}
The bounded output need not have corners. The same polynomial bound holds with a final tanh in both likelihood and posterior prediction, provided the all-width comparison class uses that same output convention.

Let $\mathcal T_{C,h}$ consist of functions $f=\tanh(g)$, where $g$ is an arbitrary-coefficient width-$C$, hidden-depth-$h$ ReLU network with a scalar affine readout. Use an unrestricted Gaussian ReLU learner with the same final tanh. All preceding posterior-certificate arguments remain valid: tanh is odd, bounded by one, and 1-Lipschitz. These three properties give, respectively, the prior mean-zero bound, the bounded-loss estimates, and the same covering and Gaussian-perturbation estimates. Its monotonicity also preserves the parameter-count bound: a threshold on $\tanh(z_w(x))$ is a threshold on $z_w(x)$ with the threshold value transformed by $\operatorname{arctanh}$. Squared-loss superlevel sets are unions of at most two such threshold sets.

\begin{theorem}[The same width polynomial with a final tanh]
\label{thm:smooth-output-oracle}
Assume $|Y|\le1$, $dP_X/dU_Q\le\kappa$, and $H\ge\max\{h,H_d\}$. Use a final tanh in both Gaussian likelihood and prediction, with the same independent Gaussian priors and two-certificate calibration. Simultaneously for all $\gamma\ge1$ and $N\ge C+1$, with probability at least $1-\delta$,
\[
 R(\widehat f_n)\le\inf_{f\in\mathcal T_{C,h}}R(f)+\varepsilon
\]
under the same sample polynomial \eqref{eq:all-width-regular-sample}, with additional logarithmic factors only.
\end{theorem}
\begin{proof}
A function $f\in\mathcal T_{C,h}$ is monotone on each affine interval of its raw network along a coordinate line. There are at most $R_0=(C+1)^h$ such intervals, so its total variation on that line is at most $2R_0$.

Fix $0<\eta<1/2$ and define
\[
 F(x)=\operatorname{arctanh}((1-\eta)f(x)),\qquad
 A_\eta=\operatorname{arctanh}(1-\eta)
       \le\tfrac12\log(2/\eta).
\]
This transformed target is bounded by $A_\eta$. It is still monotone on each of the same intervals and hence has coordinate-line total variation at most $2A_\eta R_0$. The interpolation proof only used these amplitude and variation facts. Applying the hierarchical construction to $F/A_\eta$ and multiplying its readout by $A_\eta$ gives a raw ReLU network $q$ with
\[
 \operatorname{width}(q)=O_d(\Delta^{-d}),\quad
 P(q)=O_d(A_\eta R_0\Delta^{-d}),\quad
 \int_Q|q-F|\,dU_Q\le2A_\eta dR_0\Delta.
\]
The resulting bounded output satisfies
\[
 |\tanh(q(x))-f(x)|\le|q(x)-F(x)|+\eta.
\]
Take $\eta=\varepsilon/32$ and a dyadic $\Delta$ within a factor two below $\varepsilon/[128\kappa A_\eta dR_0]$. Density domination gives a population squared-loss increase below $\varepsilon/4$. The compiler's path variation is
\[
 \widetilde O_d(\kappa^d R_0^{d+1}\varepsilon^{-d}),
\]
where the suppressed factors are powers of $A_\eta=O(\log(1/\varepsilon))$. Its width has the same polynomial bound as before, again up to logarithms.

At widths that fit this compiler, the Gaussian implicit certificate gives the same sample polynomial. At smaller widths $N\ge C+1$, the original comparison function belongs to the learner's tanh-output class, after compiling biases into a constant wire and padding identity hidden layers. The finite-parameter certificate therefore completes the same all-width argument. The two certificate events, comparison concentration, and width-confidence allocation are unchanged.
\end{proof}

There is a necessary distinction between comparison classes. A tanh-output width-$C$ network need not represent the same functions as a raw-readout width-$C$ network. The theorem above compares with networks sharing the learner's final tanh. For comparison with arbitrary raw ReLU teachers, the transformed-interpolation argument still applies after first clipping that teacher, but it only directly proves the guarantee once learner width fits the resulting compiler. The all-width theorem for raw teachers uses the hard saturation of Theorem~\ref{thm:all-width-regular-Gaussian}.

\appendix
\section{A bounded output used in likelihood and prediction}
\label{sec:saturated-specialization}
This section gives a separate learner convention. Put
\[
 c(t)=\max\{-1,\min\{t,1\}\},\qquad f_w(x)=c(z_w(x)),
\]
where $z_w$ is the ReLU network with $H$ hidden layers and total affine depth $L=H+1$, evaluated at $\xi=(1,x)/\sqrt D$, $D=d+1$. The output $f_w$, rather than $z_w$, is used in both the Gaussian likelihood and prediction. At width $N$, keep independent centered Gaussian variances $\tau^2/D$, $\tau^2/N$, and $\tau^2/N^2$ in the first, internal, and linear readout layers. Here $0<\tau\le1$ and $\omega=N/\tau^2$. For labels $|Y|\le1$, define
\[
 \Pi_{S,\lambda}(dw)=Z_S(\lambda)^{-1}
 e^{-\omega\widehat R_S(f_w)/(2\lambda)}\mu_{N,\tau}(dw),
 \qquad
 F_S(\lambda)=-\frac{2\lambda}{\omega}\log Z_S(\lambda).
\]
Thus the Gaussian observation variance is $n\lambda/\omega$. The prediction is the ordinary posterior mean of the bounded network output,
\[
 \widehat f_S(x)=\int f_w(x)d\Pi_{S,\lambda}(w).
\]
This is generally different from clipping the posterior mean of the unsaturated $z_w$.

Set the activation constant $A=1$ and define the operator-norm quantity
\[
 K(w)=\max\{128,\|W^1\|_F^2/N,\|W^2\|_{\rm op}^2,\ldots,
                 \|W^H\|_{\rm op}^2,N\|v\|^2\}.
\]
The elementary Gaussian estimates give
\[
 \mu(K>k)\le(H+1)e^{-\omega k/16},\qquad
 \log\E_\mu e^{\omega K/32}\le\omega B_H,
 \quad B_H=4+\log(H+2),
\]
for $k\ge128$. Here is a direct derivation. The first-layer quantity has law $\chi^2_{ND}/(D\omega)$, and the readout quantity has law $\chi^2_N/\omega$. Gaussian chi-square Chernoff bounds, with $N/\omega\le1$, give exponential tails at these thresholds. For an internal matrix, the Gaussian spectral-norm inequality
\[
 \Pr\{\|W^\ell\|_{\rm op}>(2\sqrt N+t)/\sqrt\omega\}
 \le2e^{-t^2/2}
\]
gives the same tail after taking $t=\sqrt{\omega k}-2\sqrt N$. A union bound and $k\ge128$ yield the displayed bound with its safe constants. Integrating this tail against $(\omega/32)e^{\omega k/32}$, with the deterministic lower cutoff $128$, gives the exponential-moment bound and the stated $B_H$.

Every raw network in the norm ball has input Lipschitz constant at most $k^{L/2}$: multiply the first Frobenius norm, the internal operator norms, and the readout norm; the factors $\sqrt N$ and $1/\sqrt N$ cancel. The fixed feature map has Lipschitz constant at most one. Its saturated output has the same Lipschitz bound and amplitude at most one. A grid on $Q$ of mesh $a/(4\sqrt d\,k^{L/2})$, with function values quantized in steps $a/2$, therefore yields the deterministic uniform-cover bound
\begin{equation}
 \log\mathcal N(a,\{f_w:K(w)\le k\},\|\cdot\|_\infty)
 \le C_{d,H,A} k^{dL/2}a^{-d}\log\frac{C_{d,H,A}k}a.
 \label{eq:saturated-cover}
\end{equation}
Nearest-grid-point extensions of the quantized values give centers in $[-1,1]$, and their squared losses are bounded by four. Conditional symmetry of the centered Gaussian readout and oddness of $c$ give $\E_\mu f_w(x)=0$. Hence
\[
 \E_\mu\widehat R_S(f_w)\le2,\quad
 \E_\mu R(f_w)\le2,\quad R(f_w)^2\le16.
\]

Put
\[
 \rho_d=(d+2)/2,\qquad
 p=\max\{d+2,dL/2\},\qquad q=p-\rho_d,\qquad r=p/\rho_d\ge2,
\]
\[
 b_n=\max\{16,4\log[(2+4\sqrt{H+1})n^3]\},\qquad
 \theta_n=C_{d,H,A}\log^6(2n/\delta)/n.
\]
The constant is fixed sufficiently large before sampling. When
$\theta_n b_n^p5^q\le1$, select the unique positive root
\begin{equation}
 \lambda_S^p=\theta_n[2F_S(\lambda_S)+\lambda_S b_n]^q.
 \label{eq:saturated-calibration}
\end{equation}
Otherwise use the deterministic fallback $\lambda_S=2\omega$.

\begin{theorem}[A saturated Gaussian posterior certificate]
\label{thm:saturated-certificate}
On one event of probability at least $1-\delta/2$, simultaneously for every $N\ge1$ and every $0<\tau\le1$, whenever the root branch is used,
\begin{equation}
 R(\widehat f_S)\le\int R(f_w)d\Pi_{S,\lambda_S}(w)
 \le F_S(\lambda_S)+2B_H\lambda_S+n^{-3}.
 \label{eq:saturated-certificate}
\end{equation}
This right side is observable from the specified prior and data. On the fallback branch the deterministic certificate $2e^{1/2}$ is valid instead. Neither branch uses teacher resources or risk.
\end{theorem}
\begin{proof}
The root argument uses only $0\le F\le2$ and monotonicity of $F/\lambda$, so it gives a unique measurable root with $1/n\le\lambda\le1/b_n$. Put $u=2F(\lambda)+\lambda b_n$; then $16\lambda\le u\le5$.

For $J=\widehat R_S(f_w)+\lambda K(w)/16$, the event $J>u$ lies in $\{\widehat R_S>u/2\}\cup\{K>8u/\lambda\}$. The prior risk bound handles the first set and Cauchy--Schwarz with $R^2\le16$ and the Gaussian norm tail handles the second. After dividing by the posterior normalizer, their total raw bounded-output risk is at most
\[
 (2+4\sqrt{H+1})\exp[\omega F/(2\lambda)-\omega u/(4\lambda)]
 \le n^{-3}.
\]

For completeness, the statistical exponent changes because the loss range is now constant. Prebuild dyadic norm shells $k$ through $160n$ and dyadic cover accuracies through a sufficiently small constant times $1/n$. Apply finite-class Bernstein to every deterministic cover in \eqref{eq:saturated-cover}, allocating total failure $\delta/2$. On $J\le u$, choose $K\le k\le2K\le32u/\lambda$ and an accuracy within a factor two of $a_*=c\lambda k/\sqrt u$. All such choices lie in the prebuilt finite list, for a sufficiently small fixed $c>0$.

The cover center has empirical risk at most a constant times $u$. Its leading Bernstein deviation and the population/empirical root-risk transfer are at most $\lambda k/32$ provided
\[
 n\ge C_{d,H,A}\log^6(2n/\delta)\,
 u^{\rho_d}\lambda^{-2\rho_d}k^{dL/2-2\rho_d}.
\]
The linear Bernstein term is absorbed using $\lambda k\le32u$. If $dL/2\ge2\rho_d$, maximize at $k\le32u/\lambda$; otherwise use $k\ge128$. The resulting condition is exactly
$n\ge C_{d,H,A}\log^6(2n/\delta)u^q\lambda^{-p}$, which follows from \eqref{eq:saturated-calibration}. Thus every network in the good-objective region satisfies
$R(f_w)\le\widehat R_S(f_w)+\lambda K(w)/16$ on the same event for all widths and scales.

Integrate this inequality, extend its nonnegative empirical-plus-norm terms to all parameters, and add the tail. Gaussian entropy duality gives
\[
 \lambda\E_\Pi K/16\le(2\lambda/\omega)\operatorname{KL}(\Pi\Vert\mu)+2B_H\lambda.
\]
The variational identity for $F$ proves \eqref{eq:saturated-certificate}, and Jensen gives the posterior-mean claim. On the fallback branch, prior risk at most two gives $Z\ge e^{-1/2}$ and posterior-average risk at most $2e^{1/2}$.
\end{proof}

\begin{corollary}[Comparison with bounded outputs of small networks]
\label{cor:saturated-oracle}
Let $g$ be any fixed bounded-output comparison network. Suppose a deterministic resource $A_{g,n}\ge1$ satisfies, for every data set and every $\lambda\ge1/n$,
\[
 F_S(\lambda)\le\widehat R_S(g)+\lambda A_{g,n}
\]
throughout the eligible widths and prior scales. The selective Gaussian measure of Section~\ref{sec:selective-standalone} supplies exactly this comparison. Define
\[
 \mathcal B_{g,n}=(A_{g,n}+b_n+2B_H+1)[1+\log(8/\delta)].
\]
Then, on one event of probability at least $1-\delta$, simultaneously for every eligible width and scale,
\begin{equation}
 R(\widehat f_S)\le R(g)+C_{d,H,A}\left[
 \mathcal B_{g,n}R(g)^{1-1/r}\theta_n^{1/p}
 +\mathcal B_{g,n}^r\theta_n^{1/\rho_d}\right]+n^{-3}.
 \label{eq:saturated-oracle}
\end{equation}
The event is uniform over every width and scale for which the deterministic comparison holds. At bounded comparison risk, a sufficient sample count for excess $\varepsilon\in(0,1)$ is
$\widetilde O((\mathcal B_{g,n}/\varepsilon)^p)$.
\end{corollary}
\begin{proof}
The proof uses the stated free-energy comparison.

The root equation, split according to which of $2\widehat R(g)$ and $(2A_{g,n}+b_n)\lambda$ is larger, yields
\[
 \lambda\le C_{d,H,A}\left[
 \theta_n^{1/p}\widehat R(g)^{1-1/r}
 +\theta_n^{1/\rho_d}(2A_{g,n}+b_n)^{r-1}\right].
\]
Substitute in the certificate. Bernstein for the bounded teacher loss, followed by
$\sqrt{RD}\le R^{1-1/r}D^{1/r}+D$ with $D=\mathcal B_{g,n}^r\theta_n^{1/\rho_d}$, gives \eqref{eq:saturated-oracle} with leading coefficient one. To see the absorption explicitly, Bernstein gives $\widehat R(g)\le R(g)+c\sqrt{R(g)D}+cD$, because $D$ dominates the scalar concentration term. The displayed inequality absorbs its square root into the two oracle remainders. Also $\widehat R(g)\le2R(g)+cD$, so raising it to $1-1/r$ and multiplying by $D^{1/r}$ gives only those same two remainders. The leading coefficient on $R(g)$ remains one.

If the root branch's condition fails, $\mathcal B_{g,n}^r\theta_n^{1/\rho_d}>5^{-q/\rho_d}$, so enlarging the fixed constant absorbs the fallback certificate. The sample claim follows because $p=r\rho_d\ge2\rho_d$: $n$ of order $(\mathcal B/\varepsilon)^p$ controls the first term and also the second, whose requirement is only $n$ of order $\mathcal B^p\varepsilon^{-\rho_d}$, up to logarithms.
\end{proof}

\section{A selective Gaussian comparison at every prior scale}
\label{sec:selective-standalone}
This section proves the free-energy comparison used above. The input is $\xi=(1,x)/\sqrt D$ with $D=d+1$, and the Gaussian precision is $\omega=N^\gamma$, $\gamma\ge1$.

\begin{lemma}[Gaussian perturbation of a sparse implementation]
Let a mean ReLU network $w_*$ have at most $m\le N$ active coordinates in every hidden layer and weighted energy $Q_N(w_*)\le E$. For $0<r\le1$, center a product Gaussian at $w_*$, retain the prior variances on coordinates entering inactive hidden rows, and multiply the prior standard deviations by $r$ on all coordinates entering active hidden rows and on every readout coordinate. Denote this comparison law by $\nu_r$. Then
\begin{align}
 \frac2\omega\operatorname{KL}(\nu_r\Vert\mu)
 &\le E+2[D+1+(H-1)m]\log(1/r),\label{eq:selective-entropy}\\
 \sup_{x\in Q}\left(\mathbb E_{\nu_r}|z_w(x)-z_{w_*}(x)|^2\right)^{1/2}
 &\le rK_E,\qquad K_E=\sqrt2(H+1)(E+1)^{H/2}.
 \label{eq:selective-rms}
\end{align}
The inactive mean rows and columns are zero, as is the mean readout off the final active coordinates.
\end{lemma}
\begin{proof}
There are at most $mD+(H-1)mN+N$ narrowed Gaussian coordinates. Since $N/\omega\le1$ and $m\le N$, their count divided by $\omega$ is at most $D+1+(H-1)m$. The Gaussian relative-entropy identity, with $r^2-1\le0$, gives \eqref{eq:selective-entropy} exactly.

Fix an input. Put $S_\ell=\|h^\ell\|_2^2/N$. The normalized first mean squared norm is at most $E$, while the first-layer noise contributes at most one to its expectation. Later mean matrices have operator norm at most $\sqrt E$. Conditional on the preceding layer, each internal noise matrix contributes at most $(N/\omega)S_{\ell-1}\le S_{\ell-1}$ to normalized expected squared norm. Centering removes cross terms and $\operatorname{ReLU}(t)^2\le t^2$, so
\[
 \mathbb E S_\ell\le2(E+1)^\ell.
\]

Let $P_\ell$ project to the active coordinates, and let
$\Delta_\ell=(\mathbb E\|P_\ell h^\ell-h_*^\ell\|_2^2/N)^{1/2}$.
The mean state is supported on those coordinates. Current-layer independence, the narrowed active-row variances, and Lipschitzness give
\[
 \Delta_1\le\sqrt2r,\qquad
 \Delta_\ell\le\sqrt E\,\Delta_{\ell-1}
       +r\sqrt{m/\omega}\sqrt{\mathbb E S_{\ell-1}}
 \le\sqrt E\,\Delta_{\ell-1}
       +\sqrt2r(E+1)^{(\ell-1)/2}.
\]
Therefore $\Delta_\ell\le\sqrt2\ell r(E+1)^{(\ell-1)/2}$. The mean readout has $\sqrt N\|v_*\|_2\le\sqrt E$; its independent perturbation contributes output variance $(r^2/\omega)S_H$. The output RMS error is thus at most
\[
 \sqrt E\,\Delta_H+r\sqrt{\mathbb E S_H/\omega}
 \le\sqrt2(H+1)r(E+1)^{H/2},
\]
proving \eqref{eq:selective-rms}.
\end{proof}

Apply the lemma to the mean implementation in \eqref{eq:path-to-Gaussian-energy}. Let $g=s(z_{w_*})$ and choose $r=(8nK_E)^{-1}$. Saturation is 1-Lipschitz, and squared loss between predictions in $[-1,1]$ and labels in $[-1,1]$ changes by at most four times their absolute prediction difference. Therefore
\[
 \mathbb E_{\nu_r}\widehat R_S(s(z_w))
 \le\widehat R_S(g)+1/(2n).
\]
The variational formula gives, for every $\lambda\ge1/n$,
\begin{equation}
 F_S(\lambda)\le\widehat R_S(g)+\lambda A,
 \qquad
 A=E+1+2[D+1+(H-1)m]\log(8nK_E).
 \label{eq:explicit-selective-cost}
\end{equation}
In particular this proves the declared resource $\widetilde O_{d,H}(1+P+m)$ without a teacher coefficient cap and without dependence on learner width.

\begin{thebibliography}{99}
\bibitem{vapnik} Vladimir N. Vapnik. \emph{Statistical Learning Theory}. Wiley, 1998. Relative uniform-convergence bounds for VC classes.
\bibitem{bartlettvc} Peter L. Bartlett, Nick Harvey, Christopher Liaw and Abbas Mehrabian. Nearly-tight VC-dimension and pseudodimension bounds for piecewise linear neural networks. \emph{Journal of Machine Learning Research} 20(63):1--17, 2019. \url{https://jmlr.org/papers/v20/18-910.html}.
\bibitem{maass} Wolfgang Maass. Bounds for the computational power and learning complexity of analog neural nets. \emph{SIAM Journal on Computing} 26(3):708--732, 1997. \url{https://doi.org/10.1137/S0097539793256041}.
\bibitem{vardi-shamir} Gal Vardi and Ohad Shamir. Neural networks with small weights and depth-separation barriers. \emph{Advances in Neural Information Processing Systems} 33:19433--19442, 2020. \url{https://papers.nips.cc/paper/2020/file/e1fe6165cad3f7f3f57d409f78e4415f-Paper.pdf}.
\end{thebibliography}
\end{document}
